\documentclass[reqno]{amsart} 
\usepackage{amssymb,latexsym,amsmath,amscd,graphicx,setspace,amsthm,verbatim}
\usepackage[margin = 3 cm]{geometry}


\theoremstyle{plain}
\newtheorem{theorem}{Theorem}[section]
\newtheorem{proposition}{Proposition}
\newtheorem{corollary}{Corollary}
\newtheorem{lemma}{Lemma}
\newtheorem{conjecture}{Conjecture}
\newtheorem{question}{Question}
\newtheorem{problem}{Problem}
      
\theoremstyle{definition}
\newtheorem{definition}{Definition}

\newenvironment{solution}{\paragraph{\emph{Solution}.}}{\hfill$\square$}


\newenvironment{solution1}{\paragraph{\emph{Solution $1$}.}}{\hfill$\square$}
\newenvironment{solution2}{\paragraph{\emph{Solution $2$}.}}{\hfill$\square$}
\newenvironment{solution3}{\paragraph{\emph{Solution $3$}.}}{\hfill$\square$}

\begin{document} 

\title[Homework 5]{Homework 5 \\ First due date: October 1.  \\  Due date for the second submission:  October 8.}

\date{\today} 
\maketitle 



\begin{problem}
In class, we gave a third way of putting a structure of a regular 1-dimensional smooth manifold on the circle $S^{1} \subseteq \mathbb{R}^{2}$ via the stereographic projection $\varphi_{N}: \mathbb{R} \rightarrow V_{N} \cap S^{1}$ from the north pole and the stereographic projection $\varphi_{S}: \mathbb{R} \rightarrow V_{S} \cap S^{1}$ from the south pole (recall that $V_{N}$ is the open set in $\mathbb{R}^{2}$ obtained by removing the non-negative part of the $y$-axis and $V_{S}$ is the open set in $\mathbb{R}^{2}$ obtained by removing the non-positive part of the $y$-axis).  I even wrote down carefully the parametrizations on the board for both $\varphi_{N}$ and $\varphi_{S}$.
\begin{enumerate}
\item Think about the picture we drew on the board and derive the formulas for both $\varphi_{N}$ and $\varphi_{S}$.
\item Then show carefully that both $\varphi_{N}$ and $\varphi_{S}$ satisfy the three required conditions in the definition of a regular manifold.
\end{enumerate}
\end{problem}
\begin{solution}

\end{solution}

\begin{problem}
Let's go back to the first way of covering the circle $S^{1}$ with four local parametrizations.  Do the same for the sphere
$$S^{2} = \{(x,y,z) \in \mathbb{R}^{3}: x^{2} + y^{2} + z^{2} = 1 \} $$
in $\mathbb{R}^{3}$.
\begin{enumerate}
\item You should get six local parametrizations whose image will be half-spheres.  Write down the six local parametrizations carefully paying close attention to the domains and codomains of your local parametrizations.  Each of them should have two parameters, since $S^{2}$ is now a surface not a curve.
\item Pick your favorite one of the six and prove carefully that the three conditions in the definition of a regular manifold are satisfied.
\item Congratulations.  You now know that $S^{2}$ is a regular 2-dimensional smooth manifold in $\mathbb{R}^{3}$.
\end{enumerate}
\end{problem}
\begin{solution}

\end{solution}


\begin{problem}
Consider the following two functions $\varphi:(0,2\pi) \times \mathbb{R} \rightarrow \mathbb{R}^{3}$ given by 
$$(\theta,t) \mapsto \varphi(\theta,t) = (\cos(\theta),\sin(\theta),t) $$
and $\psi:(-\pi,\pi) \times \mathbb{R} \rightarrow \mathbb{R}^{3}$ given by
$$(\theta,t) \mapsto \psi(\theta,t) = (\cos(\theta),\sin(\theta),t). $$
These two functions can be thought of as giving a structure of a regular $2$-dimensional smooth manifold to a certain subset $C$ of $\mathbb{R}^{3}$.
\begin{enumerate}
\item What is this subset $C$?
\item Express the codomains of both functions carefully as $V \cap C$ for some open set $V$.
\item Pick your favorite of the two functions and show that it satisfies the three conditions in the definition of a regular manifold.
\end{enumerate}
\end{problem}
\begin{proof}

\end{proof}


\begin{problem}
Consider the usual spherical parametrization of the unit sphere $S^{2}$ in $\mathbb{R}^{3}$ given by
$$(\theta,\varphi) \mapsto (\sin(\varphi) \cos(\theta), \sin(\varphi) \sin(\theta), \cos(\varphi)). $$
By modifying the domain of the function above accordingly, show that $S^{2}$ can be covered by two local parametrizations satisfying the three conditions in the definition of a regular manifold.
\end{problem}
\begin{proof}

\end{proof}














\end{document} 



